\chapter{The Rosetta Stone: Behaviour}
\label{ch:rosetta2}

\chapabstract{The second half of the reference core. Everything that
\emph{describes behaviour}: concurrent assignments, processes, sequential
statements, generate, instantiation, subprograms, packages, attributes,
reporting and time. It ends with the complete \VHDL{}-to-\SV{} mapping table
for both halves of the Rosetta Stone, which is the page to photocopy.}

\Cref{ch:rosetta} dealt with the static half of the language. Four of its
results are assumed on every page of this one, so they are worth restating
before anything moves.

\begin{itemize}
  \item \code{logic} is a single-driver four-state variable, the counterpart of
        \code{std\_ulogic}; \kw{wire} is the resolved net you need only for a
        genuine multi-driver bus.
  \item Widths never have to match. An assignment, a port connection or an
        operand of the wrong width is truncated or extended in silence, and the
        width of an arithmetic expression is decided partly by the thing you
        assign it to (\cref{sec:rosetta:widths}).
  \item Signedness is a property of an expression, not of a variable alone. One
        unsigned operand makes the whole expression unsigned.
  \item \code{'0} is the fill literal that replaces \code{(others => '0')}, and
        it is not the same as \code{1'b0}.
\end{itemize}

Everything in this chapter is about the other axis: what happens over time,
and in what order. The layout is the same as before, a \emph{rosetta} block
with \VHDL{} on the left and \SV{} on the right, then prose that says what is
the same, what is different, and what is missing.

The traps in this half are the more expensive ones, because a type error stops
the analyser and a behavioural error does not stop anything at all: a latch you
did not mean to infer, a race between two \code{always} blocks, a function
whose local variable is shared between two callers. \Cref{ch:philosophy}
explained why the language works this way; here it is as a set of
translations.


% ===========================================================================
\section{Concurrent statements}
\label{sec:rosetta:concurrent}
% ===========================================================================

Everything between \kw{begin} and \kw{end architecture} in \VHDL{} is
concurrent. Everything between the port list and \kw{endmodule} in \SV{} is
too. The set of concurrent statements is almost the same.

\subsection{Simple, conditional and selected assignment}

\begin{rosetta}
\begin{rleft}
-- simple
y <= a and b;

-- conditional (when/else)
y <= a when sel = '1' else b;

z <= "00" when op = "00" else
     "01" when op = "01" else
     "10" when op = "10" else
     "11";

-- selected (with/select)
with op select z <=
  "00" when "00",
  "01" when "01",
  "10" when "10",
  "11" when others;
\end{rleft}
\begin{rright}
// simple
assign y = a & b;

// conditional
assign y = sel ? a : b;

assign z = (op == 2'b00) ? 2'b00 :
           (op == 2'b01) ? 2'b01 :
           (op == 2'b10) ? 2'b10 :
                           2'b11;

// selected: no concurrent form.
// Use always_comb with case:
always_comb begin
  unique case (op)
    2'b00: z = 2'b00;
    2'b01: z = 2'b01;
    2'b10: z = 2'b10;
    default: z = 2'b11;
  endcase
end
\end{rright}
\end{rosetta}

\index{continuous assignment}

\VHDL{}'s simple signal assignment becomes \SV{}'s \kw{assign}, the
\emph{continuous assignment}. A continuous assignment drives its target
whenever any input changes, exactly like a concurrent signal assignment. The
target must be a net or a variable with no other driver, which matches
\VHDL{}'s single-driver rule for unresolved signals.

The conditional form maps to the ternary operator. Chained ternaries become
unreadable quickly; past two levels, write an \code{always\_comb} with an
\code{if}/\code{else if} chain. The selected form has no concurrent equivalent
at all, and its natural translation is \code{always\_comb} with a \code{case}.

Note a good \VHDL{} habit that survives: \code{with/select} forces you to cover
every choice, and the translated \code{case} should keep that discipline. Use
\code{unique case} plus a \code{default}, or a \code{default} alone. Never
write a \code{case} inside \code{always\_comb} without a \code{default} unless
every possible value of the selector is listed; the missing branch infers a
latch.

\VHDL{}-2008 allows \code{when/else} and \code{with/select} inside a process
as sequential statements. \SV{} has always allowed the ternary operator
anywhere an expression is allowed, so this needs no translation.

\subsection{Delays}

\begin{rosetta}
\begin{rleft}
-- inertial (the default): pulses
-- shorter than the delay are
-- swallowed
y <= a after 3 ns;

-- transport: every change is
-- propagated
y <= transport a after 3 ns;

-- explicit pulse rejection limit
y <= reject 1 ns inertial a
       after 3 ns;
\end{rleft}
\begin{rright}
// continuous assignment delay:
// inertial-like
assign #3 y = a;

// non-blocking assign with an
// intra-assignment delay:
// transport-like
always @(a) y <= #3 a;

// blocking, delay before the
// assignment: NOT a delay model,
// the process is simply blocked
always @(a) #3 y = a;

// gate/net delay:
wire #3 y;

// SystemVerilog has no keywords
// for transport vs inertial.
\end{rright}
\end{rosetta}

\index{delay!inertial}\index{delay!transport}

\VHDL{} distinguishes inertial and transport delay explicitly, and the default
for a signal assignment is inertial: a pulse narrower than the delay does not
propagate. \SV{} has no keywords for the distinction. A delay on a continuous
assignment behaves inertially: a new value arriving before the old one has
propagated cancels it. The closest thing to transport delay is an
\emph{intra-assignment} delay on a non-blocking assignment, \code{y <= \#3 a;}:
the right-hand side is sampled immediately and the update is scheduled three
time units later, so several pulses can be in flight at once and none is
swallowed.

Do not confuse that with \code{\#3 y = a;}, where the delay comes before the
assignment. That form models no delay at all. The process stops for three time
units, and every change of \code{a} during those three units is missed, because
the process is not back at its event control to see them. It is neither
inertial nor transport: it is a blocked process.

None of this matters for synthesisable code. Synthesis ignores every delay, and
delays in \RTL{} are a mistake, because they make simulation and synthesis
disagree. The only place a delay belongs is in a testbench, where the
distinction almost never arises. Learn the syntax so you can read gate-level
netlists produced by a synthesis tool, and do not write it.

\subsection{Blocks and named scopes}

\begin{rosetta}
\begin{rleft}
mult_stage : block
  signal partial : unsigned(15 downto 0);
begin
  partial <= a * b;
  p <= std_logic_vector(partial);
end block mult_stage;

-- a guarded block (obsolete):
-- b : block (guard_expr) begin
\end{rleft}
\begin{rright}
// named begin/end block:
begin : mult_stage
  // declarations allowed here
end : mult_stage

// but at module level you want
// a generate block:
if (1) begin : mult_stage
  logic [15:0] partial;
  assign partial = a * b;
  assign p       = partial;
end
\end{rright}
\end{rosetta}

A \VHDL{} \kw{block} groups concurrent statements and gives them a local
declarative region and a name in the hierarchy. \SV{}'s nearest equivalent at
module level is a \code{generate} block, which is what
\cref{sec:rosetta:generate} is about. Inside a procedural block, a named
\code{begin : label ... end : label} gives you a local declarative region in
the same way. Guarded blocks are obsolete in \VHDL{} too and have no \SV{}
counterpart.

% ===========================================================================
\section{Processes and procedural blocks}
\label{sec:rosetta:processes}
% ===========================================================================

\subsection{The combinational process}

\begin{rosetta}
\begin{rleft}
-- classic: explicit list
comb : process (a, b, sel) is
begin
  if sel = '1' then
    y <= a;
  else
    y <= b;
  end if;
end process comb;

-- VHDL-2008:
comb2 : process (all) is
begin
  if sel = '1' then y <= a;
  else              y <= b;
  end if;
end process comb2;
\end{rleft}
\begin{rright}
// Verilog-2001 explicit list
always @(a or b or sel) begin
  if (sel) y = a;
  else     y = b;
end

// Verilog-2001 implicit list
always @* begin ... end

// SystemVerilog, the only form
// you should write:
always_comb begin
  if (sel) y = a;
  else     y = b;
end
\end{rright}
\end{rosetta}

\index{always\_comb}\index{process!combinational}

\code{always\_comb} is the direct replacement for \code{process (all)}, and it
is better than \code{process (all)} in the ways that matter. Beyond computing
the sensitivity list for you, it carries four extra guarantees:

\begin{enumerate}
  \item It executes once at time zero, after all \code{initial} and
        \code{always} blocks have started, whether or not anything in its
        sensitivity set has changed. A plain \code{always @*} does not, so at
        time zero its outputs are \code{x} until an input moves. A \VHDL{}
        process also runs once at time zero, so this restores the \VHDL{}
        behaviour that \code{always @*} lost.
  \item Its sensitivity set includes signals read inside functions it calls.
        \code{always @*} does not follow into functions, which is a classic
        source of missing-sensitivity bugs. \VHDL{}-2008's \code{process (all)}
        has the same extended rule as \code{always\_comb}.
  \item It forbids blocking timing controls. You cannot write \code{\#5},
        \code{@(posedge clk)} or \code{wait} inside it. That turns ``I
        accidentally put a delay in my combinational logic'' into a compile
        error.
  \item Any variable written inside an \code{always\_comb} may not be written
        anywhere else in the design. This is checked at elaboration and it is
        the \SV{} equivalent of \VHDL{}'s single-driver rule for signals. It is
        the strongest reason to use \code{always\_comb} rather than
        \code{always @*}.
\end{enumerate}

\begin{gotcha}{\code{always\_comb} is not exactly \code{process (all)}}
Two differences catch people.

First, \code{always\_comb} triggers on a change in the value of anything it
reads, including through a function call, but \emph{not} on a change to
something read only through a hierarchical reference or through a
\code{ref}-passed argument. \code{process (all)} has the same limitation, so in
practice this is a wash. Do not put hierarchical references in combinational
logic.

Second, and this is the real one: in \VHDL{} a process assigns to
\emph{signals} using \code{<=}, so all the assignments in the process take
effect together at the end of the delta cycle, and reading a signal you assigned
earlier in the same process gives you the old value. In \SV{},
\code{always\_comb} should use \emph{blocking} assignments \code{=}, which take
effect immediately, so reading back a variable you assigned earlier in the same
block gives you the new value. The two models genuinely differ:
\begin{rosetta}
\begin{rleft}
process (all) is
begin
  t <= a;      -- signal
  y <= t;      -- OLD t!
end process;
-- y gets the previous value of t:
-- this infers a register, or
-- oscillates.
\end{rleft}
\begin{rright}
always_comb begin
  t = a;       // variable
  y = t;       // NEW t
end
// y = a. What you meant.
\end{rright}
\end{rosetta}
The \VHDL{} idiom for the right-hand behaviour is a \emph{variable} inside the
process, declared in its declarative part and assigned with \code{:=}. A
\SV{} \code{always\_comb} variable behaves like a \VHDL{} process variable, not
like a \VHDL{} signal. Keep that mapping in your head:
\emph{signal / \code{<=} / non-blocking} on one side, \emph{variable /
\code{:=} / blocking} on the other.
\end{gotcha}

\subsection{The clocked process}

\begin{rosetta}
\begin{rleft}
-- synchronous reset
sync : process (clk) is
begin
  if rising_edge(clk) then
    if rst = '1' then
      q <= (others => '0');
    else
      q <= d;
    end if;
  end if;
end process sync;

-- asynchronous reset
async : process (clk, rst_n) is
begin
  if rst_n = '0' then
    q <= (others => '0');
  elsif rising_edge(clk) then
    q <= d;
  end if;
end process async;
\end{rleft}
\begin{rright}
// synchronous reset
always_ff @(posedge clk) begin
  if (rst) q <= '0;
  else     q <= d;
end



// asynchronous reset
always_ff @(posedge clk or
            negedge rst_n) begin
  if (!rst_n) q <= '0;
  else        q <= d;
end
\end{rright}
\end{rosetta}

\index{always\_ff}\index{reset!asynchronous}

The structural correspondence is exact. \code{rising\_edge(clk)} becomes
\code{@(posedge clk)}; \code{falling\_edge} becomes \code{negedge}. The reset
sensitivity is in the same place in both languages: an asynchronous reset
appears in the sensitivity list, a synchronous one does not.

Two points of care.

\textbf{Polarity in the sensitivity list must match the test.} If the reset is
active low, you write \code{negedge rst\_n} in the list and \code{if (!rst\_n)}
in the body. Writing \code{posedge rst\_n} with \code{if (!rst\_n)} is legal and
produces a flip-flop that resets on the wrong edge, or, more often, a synthesis
error you will not understand.

\textbf{\code{always\_ff} restricts you usefully.} Like \code{always\_comb}, it
enforces the single-driver rule at elaboration, and it requires exactly one
event control at the top of the block. It also, by convention supported by
every tool, means ``this is a flip-flop'': lint tools will flag a blocking
assignment inside it.

\begin{gotcha}{Blocking assignment in a clocked block}
\begin{svcode}
always_ff @(posedge clk) begin
  a = d;        // blocking: a takes the new value IMMEDIATELY
  b = a;        // so b gets the NEW a: this is one flop, not two
end
always_ff @(posedge clk) begin
  a <= d;       // non-blocking: both updates happen at the end
  b <= a;       // b gets the OLD a: a two-stage shift register
end
\end{svcode}
The \VHDL{} equivalent of the second block is the familiar signal behaviour,
and it is what you almost always want. The first block does not just produce
different hardware: because the two \code{always} blocks in a design execute in
an unspecified order within the same time step, blocking assignments to
signals read by another clocked block produce a genuine race, and two
simulators can give different answers for the same code.

The rule is absolute and has no exceptions worth learning:
\begin{ruleofthumb}
Non-blocking \code{<=} in every \code{always\_ff}. Blocking \code{=} in every
\code{always\_comb}. Never mix the two in one block, and never assign the same
variable from two blocks.
\end{ruleofthumb}
\end{gotcha}

\subsection{Latches}

\code{always\_latch} exists and means what it says: it has
\code{always\_comb}'s sensitivity rules but does not require every path to
assign every output. Use it when you genuinely want a latch, which in
synchronous design is approximately never, and rely on the tool to complain
when a plain \code{always\_comb} accidentally infers one. \VHDL{} has no
equivalent keyword; an incomplete \code{if} in a combinational process infers
a latch there in exactly the same way and with exactly as little warning.

\subsection{Waiting: \code{wait}, \code{@} and \code{\#}}

\begin{rosetta}
\begin{rleft}
wait for 10 ns;
wait until rst_n = '1';
wait on a, b;
wait;                 -- forever

-- a process with no sensitivity
-- list must end in a wait
stim : process is
begin
  rst_n <= '0';
  wait for 100 ns;
  rst_n <= '1';
  wait;
end process stim;
\end{rleft}
\begin{rright}
// initial: runs once at time 0
// and then stops. No sensitivity
// list, no trailing wait needed.
initial begin
  rst_n = 1'b0;
  wait (cfg_done);   // no edge
  @(a or b);         // any change
  rst_n = 1'b1;  #100ns;
  $finish;
end
\end{rright}
\end{rosetta}

\index{wait}\index{initial}

The mapping of the waiting forms is in \cref{tab:rosetta:waits}.

\begin{table}[htbp]
  \centering
  \caption{Waiting.}
  \label{tab:rosetta:waits}
  \small
  \begin{tabular}{@{}>{\raggedright\arraybackslash}p{0.28\linewidth}>{\raggedright\arraybackslash}p{0.19\linewidth}>{\raggedright\arraybackslash}p{0.47\linewidth}@{}}
    \toprule
    \VHDL{} & \SV{} & Meaning \\
    \midrule
    \code{wait for 10 ns}      & \code{\#10ns}          & Suspend for a fixed time. \\
    \code{wait on a, b}        & \code{@(a or b)}       & Suspend until any listed signal changes. \\
    \code{wait until c = '1'}  & \code{wait (c)}        & Return immediately if already true; otherwise wait for it. \\
    \code{wait until c = '1'}  & \code{@(posedge c)}    & Not the same: this always waits for an edge. \\
    \code{wait}                & \code{wait (0)}        & Suspend forever. Rare in \SV{}. \\
    \code{wait until rising\_edge(clk)} & \code{@(posedge clk)} & The standard testbench step. \\
    ---                        & \code{wait fork}       & Wait for all processes spawned by \code{fork} in this scope. \\
    ---                        & \code{disable fork}    & Kill them instead. \\
    ---                        & \code{final}           & A block that runs once at the end of simulation. \\
    \bottomrule
  \end{tabular}
\end{table}

The subtle row is \code{wait until}. \VHDL{}'s \code{wait until c = '1'} waits
for an \emph{event} on \code{c} that makes the condition true, so if \code{c}
is already \code{'1'} the process still blocks. \SV{}'s \code{wait (c)} checks
first and returns immediately if the condition already holds. To reproduce the
\VHDL{} behaviour exactly you need \code{@(posedge c)}. Getting this wrong in a
reset sequence gives you a testbench that hangs, or one that does not wait at
all.

\code{initial} has no \VHDL{} keyword but corresponds precisely to a process
with no sensitivity list that ends in a \code{wait}. \code{final} has no
\VHDL{} equivalent at all: it is a block that runs once, at the end of
simulation, after the last event, and it is the right place for a coverage
summary or a final scoreboard check.

The last two rows before \code{final} are listed for completeness only.
\code{fork} spawns concurrent processes at run time, which a \VHDL{} process
cannot do, and \code{wait fork} and \code{disable fork} are how you join or
kill them. There is no translation to give, because there is no \VHDL{}
construct on the left; \refverificationbook introduces the whole mechanism.

\begin{gotcha}{A module-level variable is shared by every process}
A \VHDL{} process variable is declared in the process's declarative part and is
visible only there. Two processes cannot share one by accident.
\begin{rosetta}
\begin{rleft}
p1 : process is
  variable count : integer := 0;
begin
  count := count + 1;
  wait for 10 ns;
end process p1;

-- count belongs to p1.
-- p2 cannot see it. Sharing
-- requires a shared variable
-- of a protected type.
\end{rleft}
\begin{rright}
int count;    // module scope,
              // STATIC lifetime

initial forever begin
  count = count + 1; #10ns;
end
initial forever begin
  count = count + 1; #10ns;
end
// Both processes write the same
// object. Race. No warning.
\end{rright}
\end{rosetta}
Every variable declared in a module has static lifetime and module scope, so
``local to this process'' is not the default: it is something you must arrange, by
declaring the variable inside a named \code{begin}/\code{end} block, or inside
an \kw{automatic} task or function. \code{always\_comb} and \code{always\_ff}
protect you for synthesisable code by making a second writer an elaboration
error, but plain \code{always} and \code{initial} blocks do not.
\end{gotcha}

% ===========================================================================
\section{Sequential statements}
\label{sec:rosetta:sequential}
% ===========================================================================

\subsection{Conditionals}

\begin{rosetta}
\begin{rleft}
if a = '1' then
  y <= b;
elsif c = '1' then
  y <= d;
else
  y <= e;
end if;
\end{rleft}
\begin{rright}
if (a)
  y = b;
else if (c)
  y = d;
else
  y = e;
\end{rright}
\end{rosetta}

Three cosmetic differences and no semantic one: parentheses are required around
the condition, \code{elsif} is spelt \code{else if}, and there is no
\code{end if} because a multi-statement branch is delimited by
\code{begin}/\code{end}. The last is the one that produces bugs, because
forgetting \code{begin}/\code{end} silently attaches only the first statement
to the branch. Always write \code{begin}/\code{end}, even for a single
statement, and your \VHDL{} instinct will be satisfied.

\subsection{Case}

\begin{rosetta}
\begin{rleft}
case op is
  when "00" =>
    y <= a;
  when "01" | "10" =>
    y <= b;
  when others =>
    y <= c;
end case;

-- VHDL-2008 matching case:
case? addr is
  when "1---" => sel <= '1';
  when others => sel <= '0';
end case?;
\end{rleft}
\begin{rright}
case (op)
  2'b00:
    y = a;
  2'b01, 2'b10:
    y = b;
  default:
    y = c;
endcase

// wildcard case:
casez (addr)
  4'b1???: sel = 1'b1;
  default: sel = 1'b0;
endcase
\end{rright}
\end{rosetta}

\index{case!casez}\index{case!casex}

The plain \code{case} matches. The important difference is completeness:
\VHDL{} requires a \code{case} to cover every value of the selector's subtype,
so \code{when others} is compulsory unless every value is listed. \SV{} does
not require anything, and a \code{case} without a \code{default} that fails to
match simply leaves the targets unchanged, which in an \code{always\_comb}
infers a latch. Write a \code{default} in every combinational \code{case}. There
is no substitute for the discipline \VHDL{} enforced for you.

\code{casez} treats \code{z} (written \code{?}) in either the case expression or
a case item as a wildcard. \code{casex} treats both \code{x} and \code{z} as
wildcards. \VHDL{}-2008's \code{case?} matches \code{'-'} in a choice and is
the counterpart of \code{casez}.

\begin{gotcha}{Never write \code{casex}}
\code{casex} treats \code{x} as a don't-care \emph{in the case expression as
well as in the items}. If the selector is unknown, because a reset was missed
or an uninitialised register propagated, \code{casex} matches the first item
anyway and the simulation proceeds happily along a branch that the silicon will
not take.
\begin{svcode}
logic [3:0] sel;      // x after reset
casex (sel)
  4'b0001: y = a;     // MATCHES when sel is 4'bxxxx. Bug hidden.
  4'b0010: y = b;
  default: y = c;     // never reached
endcase
\end{svcode}
\code{casez} is safer because only \code{z} is a wildcard and \code{z} rarely
appears by accident inside a chip. Better still, use \code{case (sel) inside}
with \code{4'b0001:} style patterns, which is the \SV{}-2005 form and matches
only where you ask. Best of all, for a decoder, write \code{case (1'b1)} with
explicit conditions:
\begin{svcode}
always_comb begin
  unique case (1'b1)
    sel[0]: y = a;
    sel[1]: y = b;
    default: y = c;
  endcase
end
\end{svcode}
\end{gotcha}

\subsection{Case modifiers: \code{unique}, \code{unique0}, \code{priority}}

\begin{rosetta}
\begin{rleft}
-- VHDL has no modifiers. The
-- completeness requirement plus
-- the semantics of `case' already
-- imply parallel, full behaviour.

-- pragmas are needed only to
-- persuade synthesis:
case op is           -- pragma full_case
  when "00" => y <= a;
  when "01" => y <= b;
  when others => y <= c;
end case;
\end{rleft}
\begin{rright}
// exactly one item matches:
unique case (op) ... endcase

// at most one item matches:
unique0 case (op) ... endcase

// items are tried in order,
// first match wins:
priority case (op) ... endcase

// plain: first match wins, no
// check at all
case (op) ... endcase
\end{rright}
\end{rosetta}

\index{unique case}\index{priority case}

These modifiers do two things at once, and the second is why they are worth
using. To synthesis they say ``full case'' (\code{unique}, \code{priority}) and
``parallel case'' (\code{unique}), replacing the comment pragmas that Verilog
users have written for thirty years. To simulation they add a run-time check:
\code{unique} reports a violation if no item matches or if more than one
matches, \code{unique0} if more than one matches, \code{priority} if none
matches. That run-time check is genuinely valuable and it is something \VHDL{}
never offered.

\begin{tipbox}{Use \code{unique case} plus \code{default}}
The two are not redundant, but be clear about what each one buys.
\code{default} tells the compiler that every path assigns the outputs, so no
latch is inferred, and it gives you \VHDL{}'s \code{when others} completeness
guarantee. \code{unique} adds the overlap check: simulation complains if two
items match the same value.

What you do \emph{not} get is the no-match check. \code{unique} reports a
violation when no item matches, and a \code{default} makes the case always
match, so that half of the check is switched off by construction. If it is the
unreachable-state case you want to catch, drop the \code{default} and assign
the outputs before the \code{case}, or add an explicit
\code{default: \$error(...)}.
\begin{svcode}
always_comb begin
  unique case (state_q)
    IDLE: state_d = req ? RUN : IDLE;
    RUN:  state_d = done ? DONE : RUN;
    DONE: state_d = IDLE;
    default: state_d = IDLE;    // also covers x propagation
  endcase
end
\end{svcode}
\end{tipbox}

\subsection{Loops}

\begin{rosetta}
\begin{rleft}
for i in 0 to 7 loop
  y(i) <= a(7-i);
end loop;

for i in a'range loop
  ...
end loop;

while cnt < 10 loop
  cnt := cnt + 1;
end loop;

loop                -- infinite
  ...
  exit when done;
end loop;

next when skip;     -- continue
exit;               -- break
null;               -- do nothing
\end{rleft}
\begin{rright}
for (int i = 0; i < 8; i++)
  y[i] = a[7-i];

foreach (a[i])
  ...

while (cnt < 10)
  cnt = cnt + 1;

forever begin       // infinite
  ...
  if (done) break;
end

repeat (8) @(posedge clk);
do ... while (cond);

continue;
break;
;                   // do nothing
\end{rright}
\end{rosetta}

\index{loop!for}\index{foreach}

The mapping is direct. Four remarks.

\code{for} in \SV{} is the C form, with an initialisation, a condition and an
increment. Declare the loop variable in the header (\code{for (int i = 0; ...)})
and it is automatic and local to the loop; declare it outside and it is a
static module-level variable shared with everything else, which in a
testbench with several concurrent loops is a race.

\code{foreach} has no \VHDL{} counterpart and is the natural translation of
\code{for i in a'range loop}. It works on packed and unpacked arrays, on any
number of dimensions (\code{foreach (m[i][j])}), and it gets the bounds and the
direction right without you writing them.

\code{repeat (N)} has no \VHDL{} equivalent and is extremely useful in
testbenches: \code{repeat (10) @(posedge clk);} is the idiomatic ``wait ten
clocks''.

\code{do ... while} exists in \SV{} and not in \VHDL{}. \code{null} becomes an
empty statement, a bare semicolon, which is easy to write by accident: an
\code{if (c);} with a stray semicolon compiles and does nothing.

For synthesis, both languages require loop bounds to be static, and the loop is
unrolled. The \SV{} rule is the same as the \VHDL{} one.

% ===========================================================================
\section{Generate}
\label{sec:rosetta:generate}
% ===========================================================================

\begin{rosetta}
\begin{rleft}
gen_stages : for i in 0 to N-1 generate
  u_st : entity work.stage
    generic map (IDX => i)
    port map (
      clk => clk,
      d   => pipe(i),
      q   => pipe(i+1));
end generate gen_stages;

gen_first : if i = 0 generate
  ...
end generate gen_first;

-- VHDL-2008 adds else/elsif:
gen_c : case MODE generate
  when 0 => ...
end generate gen_c;
\end{rleft}
\begin{rright}
for (genvar i = 0; i < N; i++)
  begin : gen_stages
    stage #(.IDX(i)) u_st (
      .clk (clk),
      .d   (pipe[i]),
      .q   (pipe[i+1]));
  end

if (MODE == 0) begin : gen_a
  ...
end else begin : gen_b
  ...
end

case (MODE)
  0: begin : gen_c0 ... end
  default: begin : gen_cd ... end
endcase
\end{rright}
\end{rosetta}

\index{generate}\index{genvar}

The three forms correspond one to one, including \code{case ... generate},
which \VHDL{}-2008 added and \SV{} has always had. The keywords
\code{generate} and \code{endgenerate} are optional in \SV{}-2009 and later;
most modern code omits them, because the \code{for}, \code{if} or \code{case}
at module level is unambiguously a generate construct.

The loop variable is a \kw{genvar}\index{genvar}: an elaboration-time integer,
not a run-time object. Declare it in the loop header
(\code{for (genvar i = 0; ...)}) so its scope is the loop. \VHDL{} declares the
equivalent implicitly in the \code{for ... generate} header, which is the same
idea.

\subsection{Names matter}

\begin{gotcha}{An unnamed generate block still has a name}
In \VHDL{} the generate label is compulsory. In \SV{} it is optional, and if
you leave it out the tool invents one: \code{genblk1}, \code{genblk2}, numbered
in order of appearance in the file. Every hierarchical path, every waveform
name, every synthesis constraint and every assertion binding then depends on
that number, and inserting an unrelated \code{if} earlier in the file
renumbers everything below it.
\begin{svcode}
for (genvar i = 0; i < N; i++) begin        // becomes genblk1
  stage u_st (...);
end
// path: top.genblk1[3].u_st  -- and it changes if you edit the file above

for (genvar i = 0; i < N; i++) begin : g_stage
  stage u_st (...);
end
// path: top.g_stage[3].u_st  -- stable
\end{svcode}
Name every generate block, with a consistent prefix such as \code{g\_}. It
costs nothing and it is the difference between a constraint file that survives
a refactor and one that does not.
\end{gotcha}

\subsection{A pipeline, both ways}

The canonical use of \code{generate} is a regular structure whose stages differ
only by index. Here is a systolic multiply-accumulate chain, in both languages.

\begin{vhdlcode}[caption={A generated systolic chain in \VHDL{}.},label={lst:rosetta:systolic-vhdl}]
architecture rtl of mac_chain is
  type acc_arr_t is array (0 to N) of signed(31 downto 0);
  type dat_arr_t is array (0 to N) of signed(15 downto 0);
  signal acc : acc_arr_t;
  signal dat : dat_arr_t;
begin
  acc(0) <= (others => '0');
  dat(0) <= din;

  gen_taps : for i in 0 to N-1 generate
    tap : process (clk) is
    begin
      if rising_edge(clk) then
        dat(i+1) <= dat(i);
        acc(i+1) <= acc(i) + resize(dat(i) * coeff(i), 32);
      end if;
    end process tap;
  end generate gen_taps;

  dout <= std_logic_vector(acc(N));
end architecture rtl;
\end{vhdlcode}

\begin{svcode}[caption={The same chain in \SV{}.},label={lst:rosetta:systolic-sv}]
module mac_chain #(parameter int unsigned N = 8) (
  input  wire                clk,
  input  wire signed [15:0]  din,
  input  wire signed [15:0]  coeff [0:N-1],
  output logic signed [31:0] dout
);
  logic signed [31:0] acc [0:N];
  logic signed [15:0] dat [0:N];

  assign acc[0] = '0;
  assign dat[0] = din;

  for (genvar i = 0; i < N; i++) begin : g_taps
    always_ff @(posedge clk) begin
      dat[i+1] <= dat[i];
      acc[i+1] <= acc[i] + 32'($signed(dat[i]) * $signed(coeff[i]));
    end
  end

  assign dout = acc[N];
endmodule
\end{svcode}

The two are the same design. Note three translation details. The \VHDL{} array
types become plain unpacked arrays with no \code{type} declaration. The
\code{resize(...)} disappears because the thirty-two-bit target does the
widening, but the \code{32'(...)} size cast is written anyway, so that the
multiply is done at thirty-two bits regardless of context. And the
\code{coeff} port is an unpacked array of signed vectors, which \SV{} allows
across a port boundary and which older Verilog did not.

% ===========================================================================
\section{Instantiation}
\label{sec:rosetta:instantiation}
% ===========================================================================

\subsection{Component declaration or direct instantiation}

\begin{rosetta}
\begin{rleft}
architecture rtl of top is
  -- the old way: declare a
  -- component, then instantiate it
  component fifo is
    generic (WIDTH : positive);
    port (clk : in  std_logic;
          d   : in  std_logic_vector;
          q   : out std_logic_vector);
  end component fifo;
begin
  u_fifo : fifo
    generic map (WIDTH => 16)
    port map (clk => clk,
              d   => din,
              q   => dout);
end architecture rtl;
\end{rleft}
\begin{rright}
module top (...);

  // there is no component
  // declaration in SystemVerilog.
  // The module itself is the
  // declaration.

  fifo #(
    .WIDTH (16)
  ) u_fifo (
    .clk (clk),
    .d   (din),
    .q   (dout)
  );

endmodule
\end{rright}
\end{rosetta}

\index{instantiation}\index{component}

The \VHDL{} component declaration is a local re-declaration of an entity's
interface, and it exists because \VHDL{} originally insisted that everything be
declared before use in the current design unit. Getting the component
declaration out of step with the entity is a classic \VHDL{} bug: the two are
checked only at elaboration, and a mismatched port list produces an error a
long way from its cause.

\SV{} has nothing of the kind. You instantiate a module by name, and the tool
finds it in whatever was compiled. If it is not there, most tools issue a
warning and create a black box, which is another reason to run
\tool{Verilator}, whose default is an error.

\VHDL{}-2008's direct entity instantiation removes the component declaration on
the \VHDL{} side too, and it is what you should already be writing:

\begin{vhdlcode}[caption={Direct entity instantiation: the modern \VHDL{} form.},label={lst:rosetta:direct}]
u_fifo : entity work.fifo(rtl)      -- (rtl) selects the architecture
  generic map (WIDTH => 16)
  port map (clk => clk, d => din, q => dout);
\end{vhdlcode}

The architecture name in parentheses has no \SV{} equivalent, because a module
has only one body.

\subsection{Named, positional, and shorthand connections}

\begin{rosetta}
\begin{rleft}
-- named (always do this)
port map (clk => clk,
          d   => din,
          q   => dout,
          err => open);

-- positional (do not)
port map (clk, din, dout, open);

-- VHDL has no shorthand for
-- "connect to a signal of the
-- same name".
\end{rleft}
\begin{rright}
// named (always do this)
.clk (clk),
.d   (din),
.q   (dout),
.err ()          // unconnected

// positional (do not)
fifo u (clk, din, dout, );

// implicit .name
.clk, .d, .q     // same-name

// implicit .*
fifo u_fifo (.*);
\end{rright}
\end{rosetta}

Named association is the same idea in both languages, and it is the only form
you should use. Note the spelling of an unconnected port: \VHDL{} writes
\code{open}, \SV{} writes empty parentheses \code{.err()}. Omitting the port
entirely is also legal in \SV{} and means the same thing, which is why
\cref{sec:rosetta:ports} insisted you connect everything explicitly.

\SV{} adds two shorthands that \VHDL{} does not have.

\code{.name} connects the port \code{name} to a signal called \code{name} in
the instantiating scope. \code{.*} does that for every port not explicitly
connected. Both are elaboration-time checks: if no signal of that name exists,
or the widths disagree, it is an error, not a silent implicit net.

\code{.*} is genuinely controversial. In its favour: a hundred-port CPU core
instantiated with \code{.*} is four lines instead of a hundred, name mismatches
become errors, and adding a port to the module does not require editing every
instantiation. Against it: you cannot see from the instantiation what is
connected to what, code review becomes harder, and it encourages a naming
convention where every signal in the parent has the same name as the port,
which then breaks the moment you need two instances.

\begin{ruleofthumb}
Use full named association in \RTL{} you expect a human to read. Use \code{.*}
only in testbench wrappers, where the connection is one-to-one by construction.
\end{ruleofthumb}

\subsection{Configurations and \kw{bind}}

A \VHDL{} configuration binds component instances to entity/architecture pairs,
possibly deep in the hierarchy, without editing the design. \SV{} has no
equivalent for that job. What it does have, and \VHDL{} does not, is
\kw{bind}\index{bind}: a way of attaching a module instance to an existing
module or instance from outside it.

\begin{svcode}[caption={\code{bind}: adding checkers without touching the \RTL{}.},label={lst:rosetta:bind}]
// fifo_checker is an ordinary module full of assertions.
// This attaches one to every instance of `fifo', with access to its internals:
bind fifo fifo_checker u_chk (.clk(clk), .wr(wr_en), .rd(rd_en), .cnt(count_q));

// or to one specific instance:
bind top.u_dma.u_fifo fifo_checker u_chk (.*);
\end{svcode}

\code{bind} is how assertion and coverage code is kept out of synthesisable
files. It is one of the most useful features \SV{} has that \VHDL{} lacks
entirely. \cref{ch:beyond} treats it properly and \refverificationbook puts it to
work.

% ===========================================================================
\section{Subprograms}
\label{sec:rosetta:subprograms}
% ===========================================================================

\begin{rosetta}
\begin{rleft}
function parity (v : std_logic_vector)
  return std_logic is
  variable r : std_logic := '0';
begin
  for i in v'range loop
    r := r xor v(i);
  end loop;
  return r;
end function parity;

procedure drive (
  signal   clk : in  std_logic;
  signal   d   : out std_logic_vector;
  constant v   : in  std_logic_vector) is
begin
  wait until rising_edge(clk);
  d <= v;
end procedure drive;
\end{rleft}
\begin{rright}
function automatic logic parity (
    input logic [31:0] v);
  logic r = 1'b0;
  foreach (v[i]) r ^= v[i];
  return r;
endfunction

task automatic drive (
    ref   logic        clk,
    output logic [7:0] d,
    input logic [7:0]  v);
  @(posedge clk);
  d = v;
endtask
\end{rright}
\end{rosetta}

\index{function}\index{task}

\subsection{Function or task?}

A \VHDL{} \kw{function} returns a value and cannot wait. A \VHDL{}
\kw{procedure} does not return a value and may wait. That is exactly the \SV{}
distinction between \kw{function} and \kw{task}, with one relaxation and one
addition.

The relaxation: a \SV{} \kw{function} may have \code{output} and \code{inout}
arguments as well as a return value, so it can pass several results back
without a \kw{task}. A function that returns nothing is declared
\code{function void f(...)} and is called as a statement.

The addition: a \SV{} \kw{task} may consume time (\code{\#}, \code{@},
\code{wait}) exactly as a \VHDL{} procedure may, but a \kw{function} may not,
and a synthesisable function may not either. The rule is enforced by the
compiler.

\begin{table}[htbp]
  \centering
  \caption{Subprograms.}
  \label{tab:rosetta:subprograms}
  \small
  \begin{tabular}{@{}>{\raggedright\arraybackslash}p{0.30\linewidth}>{\raggedright\arraybackslash}p{0.28\linewidth}>{\raggedright\arraybackslash}p{0.36\linewidth}@{}}
    \toprule
    \VHDL{} & \SV{} & Note \\
    \midrule
    \code{function} & \code{function} & Cannot consume time in either. \\
    \code{procedure} & \code{task} & May consume time in both. \\
    \code{procedure} (no waits) & \code{function void} & Callable as a statement. \\
    \code{impure function} & (default) & Every \SV{} function may read module state. \\
    \code{pure function} & --- & No way to declare purity. \\
    parameter mode \code{in} & \code{input} & Copied in. \\
    parameter mode \code{out} & \code{output} & Copied out on return. \\
    parameter mode \code{inout} & \code{inout} & Copied both ways. \\
    \code{signal} parameter & \code{ref} & Passed by reference; the only way to see live changes. \\
    default argument values & default argument values & Same. \\
    overloading by signature & --- & \SV{} has no subprogram overloading. \\
    operator overloading & --- & No equivalent at all. \\
    recursion (always allowed) & \code{automatic} required & See the trap below. \\
    \bottomrule
  \end{tabular}
\end{table}

\subsection{Purity}

\VHDL{} distinguishes \kw{pure} and \kw{impure} functions: a pure function may
read only its arguments, an impure one may read signals and variables in the
enclosing scope. The default is pure, and the compiler enforces it.

\SV{} has no such distinction. Every function may read anything in scope, and
nothing tells you whether it does. Where \VHDL{} would have refused to compile
a function that reads a signal you forgot to pass, \SV{} compiles it and the
function silently depends on module state. In an \code{always\_comb} this is
usually harmless, because \code{always\_comb} extends its sensitivity into
functions, but it makes the function unusable anywhere else.

\begin{ruleofthumb}
Write every \SV{} function as if it were \VHDL{}-pure: take everything through
arguments, return everything through the return value. Then it is reusable and
testable.
\end{ruleofthumb}

\subsection{Lifetime: the one-word difference}

\begin{gotcha}{Write \kw{automatic} on every subprogram}
In \VHDL{}, a subprogram's local variables are created on entry and destroyed
on return, so two concurrent calls have two independent sets and recursion
works. In \SV{} that behaviour is opt-in. A function or task declared in a
module, interface or package without the word \kw{automatic} has \emph{static}
lifetime: its locals are allocated once, for the whole simulation, and every
call shares them.

Recursion is the immediate victim, and it is the one that shows up while you
are still writing \RTL{}:
\begin{svcode}
// WRONG: static by default, so the recursive call overwrites `n'
function int factorial (int n);
  if (n <= 1) return 1;
  return n * factorial(n - 1);
endfunction

// RIGHT
function automatic int factorial (int n);
  if (n <= 1) return 1;
  return n * factorial(n - 1);
endfunction
\end{svcode}
The other victim is concurrency, which matters once you start writing
testbenches: two processes calling the same static task share its arguments and
locals, and the failure looks like random data corruption. That is the trap in
detail, together with the enclosing-unit form
(\code{module automatic tb;}) and the reason class methods are exempt, in
\refverificationbook.

For a subprogram, the rule needs no detail at all. Write \kw{automatic} on
every function and task. There is no downside: \kw{automatic} is what \VHDL{}
always did.
\end{gotcha}

\subsection{Argument passing and \code{ref}}

A \VHDL{} \code{signal} parameter lets a procedure see a signal live and drive
it, which is how testbench procedures work. The \SV{} equivalent is \kw{ref}: a
reference argument, not a copy. \code{const ref} passes by reference without
allowing modification, which is the efficient way to pass a large array into a
function.

\begin{svcode}[caption={Argument modes.},label={lst:rosetta:argmodes}]
function automatic void stats (
    const ref logic [31:0] data [],   // by reference, read only, no copy
    output int unsigned    n_ones,    // copied out on return
    inout  int unsigned    running    // copied in and out
);
  n_ones = 0;
  foreach (data[i]) n_ones += $countones(data[i]);
  running += n_ones;
endfunction
\end{svcode}

Note that \code{output} and \code{inout} arguments of a \SV{} subprogram are
copied on entry and on return, not updated live. A \kw{task} that waits and
then writes an \code{output} argument does not update the caller's variable
until it returns. For live updates you need \kw{ref}, and \kw{ref} arguments
are not allowed in a function that is called from a continuous assignment.

\subsection{Overloading}

\VHDL{} overloads subprograms by signature and lets you overload operators:
this is how \code{numeric\_std} gives \code{"+"} a meaning for
\code{unsigned}. \SV{} has neither. Two functions in the same scope may not
share a name, and there is no way to give \code{+} a new meaning.

The consequence for a \VHDL{} designer is concrete: the library-based approach
to fixed-point arithmetic (\code{ieee.fixed\_pkg}, with \code{sfixed} and
\code{ufixed} types and overloaded operators) has no \SV{} counterpart. You
manage binary points by hand, with named functions and careful comments. This
is one of the places where \VHDL{} is simply the better language, and it is
worth saying so.

% ===========================================================================
\section{Packages and visibility}
\label{sec:rosetta:packages}
% ===========================================================================

\begin{rosetta}
\begin{rleft}
package demo_pkg is
  constant ITER : positive := 14;
  subtype data_t is
    signed(15 downto 0);
  -- deferred constant: value in
  -- the body
  constant TABLE_SIZE : positive;
  function atan_rom (i : natural)
    return data_t;
end package cordic_pkg;

package body cordic_pkg is
  constant TABLE_SIZE : positive := 16;
  function atan_rom (i : natural)
    return data_t is
  begin
    return to_signed(0, 16);
  end function atan_rom;
end package body demo_pkg;
\end{rleft}
\begin{rright}
package demo_pkg;
  localparam int unsigned ITER = 14;
  typedef logic signed [15:0] data_t;

  // no deferred constants: the
  // value must be here
  localparam int unsigned
    TABLE_SIZE = 16;

  function automatic data_t
      atan_rom (int unsigned i);
    return '0;
  endfunction
endpackage
\end{rright}
\end{rosetta}

\index{package}

\subsection{One unit, not two}

A \VHDL{} package is two design units: the declaration, which is the interface,
and the body, which holds the subprogram implementations and the values of
deferred constants. Changing the body does not force recompilation of anything
that uses the package, which is a real benefit on a large design.

A \SV{} package is one unit. There is no body and no separation between
interface and implementation. Everything, including function bodies, is in the
package, and everything that imports the package is recompiled when it changes.

Deferred constants therefore have no equivalent. A \VHDL{} package can declare
\code{constant TABLE\_SIZE : positive;} in the declaration and give it a value
in the body, hiding the value from users while still letting them use the name.
\SV{} cannot do that at all.

\subsection{Import and visibility}

\begin{rosetta}
\begin{rleft}
use work.cordic_pkg.all;
use work.cordic_pkg.data_t;

-- fully qualified, no use clause
signal x : work.cordic_pkg.data_t;

-- a use clause in an entity
-- applies to the entity and to
-- every architecture of it
\end{rleft}
\begin{rright}
import cordic_pkg::*;
import cordic_pkg::data_t;

// fully qualified, no import
cordic_pkg::data_t x;

// an import before the module
// header is visible in the port
// list too:
module m
  import cordic_pkg::*;
  (input data_t d);
\end{rright}
\end{rosetta}

The scope resolution operator is \code{::}, not \code{.}, and a fully qualified
name needs no import at all. That is the safest form for occasional use, and it
documents where the name comes from.

The visibility rules differ in a way that matters. A wildcard import
(\code{import pkg::*}) does not declare anything. It makes names
\emph{candidates}: when the compiler meets an identifier that is not declared
locally and is not visible from an enclosing scope, it looks through the
wildcard imports. If exactly one supplies the name, that name is used and
becomes permanently imported into the scope. If two supply it, that reference
is an error, but only that reference: the rest of the file compiles.

An explicit import (\code{import pkg::data\_t}) does declare the name. Two
explicit imports of the same name into one scope are an error at the
declaration, whether or not you use it. A local declaration always wins over a
wildcard import, and always conflicts with an explicit one.

\begin{gotcha}{A local name silently shadows a wildcard import}
\begin{svcode}
package pkg_a;  localparam int WIDTH = 32;  endpackage

module m;
  import pkg_a::*;
  localparam int WIDTH = 8;    // shadows pkg_a::WIDTH. No warning.
  logic [WIDTH-1:0] bus;       // 8 bits, not 32.
endmodule
\end{svcode}
\VHDL{}'s \code{use ... all} behaves the same way in that a locally declared
homograph hides the used one, but \VHDL{} makes the two candidates ambiguous
where \SV{} silently prefers the local. In a large design with several
imported packages, this is how two modules end up disagreeing about a bus
width. Prefix package constants (\code{CFG\_WIDTH}), or use the qualified form
\code{pkg\_a::WIDTH} wherever ambiguity is possible.
\end{gotcha}

\code{export pkg\_a::*;} inside a package re-exports names that the package
itself imported, so that importing the outer package makes them visible too.
\VHDL{} has no equivalent: a \code{use} clause in a package declaration does
not propagate to users of that package. Use \code{export} sparingly; it makes
the origin of a name very hard to find.

\subsection{Package or header file?}

\begin{tipbox}{When to use a package and when to use \file{.svh}}
Use a \textbf{package} for everything the language can express: types,
\code{typedef}s, constants (\kw{localparam}), functions, tasks, enumerations
and \code{covergroup} definitions. A package is compiled, type-checked, and
namespaced.

Use an \textbf{included header} (\file{.svh}) only for text the language cannot
express: macros, and the conditional-compilation guards around them. Guard it
against double inclusion:
\begin{svcode}
`ifndef CORDIC_DEFS_SVH
`define CORDIC_DEFS_SVH
`define CORDIC_ASSERT(c) assert(c) else $error("cordic: %s", `"c`")
`endif
\end{svcode}
A package cannot hold macros, because macros are gone before the compiler sees
the package; a header cannot hold a type safely, because including it twice
declares the type twice. The two mechanisms do not overlap and neither replaces
the other.
\end{tipbox}

Packages themselves may be parameterised only indirectly: \SV{} has no generic
packages. \VHDL{}-2008 does, though tool support is uneven. The usual \SV{}
workaround is a package of \code{parameter type} definitions plus a class or a
module that carries the parameters, and it is markedly less elegant than the
\VHDL{} form.

% ===========================================================================
\section{Attributes and system functions}
\label{sec:rosetta:attributes}
% ===========================================================================

\VHDL{} attributes and \SV{} system functions overlap without matching. The
mapping is in \cref{tab:rosetta:attributes}. Two general remarks first.

\VHDL{} attributes are part of the language and are evaluated by the same rules
as everything else; \code{s'event} is an expression you can use anywhere.
\SV{}'s equivalents are split between system functions usable anywhere
(\code{\$bits}, \code{\$clog2}) and sampled-value functions
(\code{\$rose}, \code{\$past}, \code{\$stable}) that are meaningful only in a
clocked assertion context or with an explicit clocking event. You cannot write
\code{if (\$rose(a))} in an ordinary \code{always\_comb} and get \VHDL{}'s
\code{a'event} behaviour.

Conversely, \SV{} has a set of bit-manipulation functions with no \VHDL{}
counterpart at all, and they are used constantly: \code{\$clog2},
\code{\$countones}, \code{\$onehot}, \code{\$isunknown}.

\begin{longtable}{@{}>{\raggedright\arraybackslash}p{0.30\linewidth}>{\raggedright\arraybackslash}p{0.25\linewidth}>{\raggedright\arraybackslash}p{0.39\linewidth}@{}}
  \caption{Attributes and system functions.}
  \label{tab:rosetta:attributes} \\
  \toprule
  \VHDL{} & \SV{} & Note \\
  \midrule
  \endfirsthead
  \multicolumn{3}{@{}l}{\small\itshape \cref{tab:rosetta:attributes}, continued}\\
  \toprule
  \VHDL{} & \SV{} & Note \\
  \midrule
  \endhead
  \midrule
  \multicolumn{3}{r@{}}{\small\itshape continues}\\
  \endfoot
  \bottomrule
  \endlastfoot
  \code{s'event}          & \code{@(s)} or \code{\$changed(s)} & No expression form outside assertions. \\
  \code{rising\_edge(s)}  & \code{\$rose(s)}      & In assertions. In \RTL{}, \code{posedge}. \\
  \code{falling\_edge(s)} & \code{\$fell(s)}      & \\
  \code{s'stable}         & \code{\$stable(s)}    & Sampled-value function; needs a clock. \\
  \code{s'last\_value}    & \code{\$past(s)}      & Not identical: \code{\$past} is the value one clock ago, \code{'last\_value} the value before the last change. \\
  \code{s'delayed(t)}     & ---                   & No equivalent. \\
  \code{s'last\_event}    & ---                   & No equivalent. \\
  \code{t'image(x)}       & \code{\$sformatf("\%0d",x)}, \code{\%p} & For enumerations, \code{x.name()}. \\
  \code{t'value(str)}     & ---                   & No equivalent; \code{\$sscanf} is the nearest. \\
  \code{t'pos(x)}         & \code{int'(x)}        & \\
  \code{t'val(i)}         & \code{t'(i)}          & Or \code{\$cast} for a checked conversion. \\
  \code{t'succ(x)}, \code{t'pred(x)} & \code{x.next()}, \code{x.prev()} & Enumerations only. \\
  \code{t'leftof}, \code{t'rightof}  & ---        & \\
  \code{a'range}          & ---                   & Use \code{foreach} or \code{\$size}. \\
  \code{a'reverse\_range} & ---                   & No equivalent. \\
  \code{a'length}         & \code{\$size(a)}      & \code{\$bits(a)} for the total bit count. \\
  \code{a'high}, \code{a'low} & \code{\$high(a)}, \code{\$low(a)} & \\
  \code{a'left}, \code{a'right} & \code{\$left(a)}, \code{\$right(a)} & \\
  \code{a'ascending}      & ---                   & Test \code{\$left(a) < \$right(a)}. \\
  ---                     & \code{\$bits(t)}      & Bit count of any type, including a struct. \\
  ---                     & \code{\$clog2(n)}     & Ceiling of $\log_2 n$. Constant expression. \\
  ---                     & \code{\$countones(v)} & Population count. \\
  ---                     & \code{\$onehot(v)}, \code{\$onehot0(v)} & Exactly one, or at most one, bit set. \\
  ---                     & \code{\$isunknown(v)} & True if any bit is \code{x} or \code{z}. \\
  ---                     & \code{\$typename(x)}  & The type of \code{x} as a string. Debug only. \\
  ---                     & \code{\$dimensions}, \code{\$unpacked\_dimensions} & Array introspection. \\
  \code{maximum}, \code{minimum} (2008) & ---     & No system function; write your own. \\
\end{longtable}

\begin{tipbox}{\code{\$clog2} replaces the \code{log2} dance}
The \VHDL{} idiom for an address width is
\code{integer(ceil(log2(real(DEPTH))))} with \code{use ieee.math\_real.all},
and it has a well-known edge case at exact powers of two. \SV{} gives you
\code{\$clog2(DEPTH)} as a constant function usable directly in a
\kw{localparam}, and \code{\$clog2(1)} is 0, which is what you want. Use it in
every parameterised memory:
\begin{svcode}
localparam int unsigned AW = (DEPTH > 1) ? $clog2(DEPTH) : 1;
logic [AW-1:0] rd_ptr, wr_ptr;
\end{svcode}
\end{tipbox}

% ===========================================================================
\section{Reporting, assertions and files}
\label{sec:rosetta:report}
% ===========================================================================

\subsection{Assertions and severity}

\begin{rosetta}
\begin{rleft}
assert cnt < DEPTH
  report "fifo overflow"
  severity error;

assert false
  report "unreachable"
  severity failure;

report "starting test"
  severity note;

-- severity levels:
--   note, warning, error, failure
\end{rleft}
\begin{rright}
assert (cnt < DEPTH)
  else $error("fifo overflow");

assert (0)
  else $fatal(1, "unreachable");

$info("starting test");

// severity tasks:
//   $info, $warning, $error,
//   $fatal
\end{rright}
\end{rosetta}

\index{assertion!immediate}

The structural correspondence is exact and the mapping of severities is in
\cref{tab:rosetta:severity}. In both languages the assertion fires when the
condition is \emph{false}. The \SV{} form puts the message in an \code{else}
clause, which reads oddly at first and is logical once you see it: the
\code{else} branch is what happens when the assertion fails, and it may contain
any statement, a message being only the usual one.

\begin{table}[htbp]
  \centering
  \caption{Severity.}
  \label{tab:rosetta:severity}
  \small
  \begin{tabular}{@{}ll>{\raggedright\arraybackslash}p{0.50\linewidth}@{}}
    \toprule
    \VHDL{} & \SV{} & Behaviour \\
    \midrule
    \code{note}    & \code{\$info}    & Message only. \\
    \code{warning} & \code{\$warning} & Message only. \\
    \code{error}   & \code{\$error}   & Message; simulation continues unless the tool is told otherwise. \\
    \code{failure} & \code{\$fatal}   & Message and immediate stop. First argument is the exit code passed to \code{\$finish}. \\
    \bottomrule
  \end{tabular}
\end{table}

Two \SV{} refinements have no \VHDL{} counterpart. A \emph{deferred} assertion,
written \code{assert \#0 (...)}, is evaluated at the end of the time step rather
than immediately, which suppresses spurious failures caused by transient values
in combinational logic. A \emph{final} assertion, \code{assert final (...)},
is checked at the end of simulation. Both are worth knowing. Everything above
is an \emph{immediate} assertion, a statement that runs when control reaches
it; \SV{} also has \emph{concurrent} assertions, which describe behaviour over
time and are a sublanguage \VHDL{} has nothing comparable to.
\refverificationbook is where that sublanguage is taught, and where
the assertions written with it are wired into a testbench.

\subsection{Printing}

\begin{rosetta}
\begin{rleft}
use std.textio.all;

report "count = " &
  integer'image(cnt);

-- VHDL-2008 to_string:
report "vec = " & to_string(v);

-- printing a real
report "f = " & real'image(f);

-- hex needs a helper
report "h = " & to_hstring(v);
\end{rleft}
\begin{rright}
$display("count = %0d", cnt);

$display("vec = %b", v);

$display("f = %f", f);

$display("h = %h", v);

$write("no newline");
msg = $sformatf("%0d", cnt);
$display("%s at %0t",
         state.name(), $time);
\end{rright}
\end{rosetta}

\index{system task}

\VHDL{} builds a string by concatenation and image functions; \SV{} uses C-like
format specifiers. \cref{tab:rosetta:format} lists the ones you need.

\begin{table}[htbp]
  \centering
  \caption{Format specifiers.}
  \label{tab:rosetta:format}
  \small
  \begin{tabular}{@{}ll>{\raggedright\arraybackslash}p{0.50\linewidth}@{}}
    \toprule
    \SV{} & \VHDL{} equivalent & Meaning \\
    \midrule
    \code{\%d} & \code{integer'image} & Decimal, padded to the width of the type. \\
    \code{\%0d} & \code{integer'image} & Decimal, no padding. Use this one. \\
    \code{\%h}, \code{\%0h} & \code{to\_hstring} (2008) & Hexadecimal. \\
    \code{\%b}, \code{\%0b} & \code{to\_string} & Binary. \\
    \code{\%o} & \code{to\_ostring} (2008) & Octal. \\
    \code{\%s} & (the string itself) & String, or an enumeration's \code{.name()}. \\
    \code{\%f}, \code{\%e}, \code{\%g} & \code{real'image} & Real. \\
    \code{\%t} & \code{time'image} & Time, formatted by \code{\$timeformat}. \\
    \code{\%p} & --- & ``Pretty print'': any aggregate, including structs and arrays. No \VHDL{} equivalent. \\
    \code{\%m} & --- & The hierarchical name of the current scope. Very useful in messages. \\
    \bottomrule
  \end{tabular}
\end{table}

\code{\%0d} rather than \code{\%d} is worth a habit: \code{\%d} pads to the
width of the type, so a 32-bit counter prints as ten characters with nine
leading spaces. \code{\%p} is the closest thing \SV{} has to a universal
\code{'image}, and it will print a whole packed structure or unpacked array in
one go, which is invaluable in a scoreboard message.

\subsection{Files}

\begin{rosetta}
\begin{rleft}
use std.textio.all;

file f      : text;
variable l  : line;
variable v  : integer;

file_open(f, "in.txt", read_mode);
while not endfile(f) loop
  readline(f, l);
  read(l, v);
end loop;
file_close(f);

file_open(f, "out.txt", write_mode);
write(l, string'("hello"));
writeline(f, l);
file_close(f);
\end{rleft}
\begin{rright}
int fd;
int v;
string s;

fd = $fopen("in.txt", "r");
while (!$feof(fd)) begin
  void'($fscanf(fd, "%d\n", v));
end
$fclose(fd);

fd = $fopen("out.txt", "w");
$fdisplay(fd, "hello");
$fclose(fd);

// memory initialisation:
$readmemh("init.hex", mem);
$readmemb("init.bin", mem);
\end{rright}
\end{rosetta}

\index{file I/O}\index{readmemh}

\code{std.textio}'s two-step model, read a \code{line} object and then parse
fields out of it, has no \SV{} counterpart. \SV{} reads directly with
\code{\$fscanf} (formatted), \code{\$fgets} (a whole line into a
\kw{string}), or \code{\$fread} (binary). The file handle is an \kw{int}, not a
distinct file type, and passing a stale or zero handle is not caught.

\code{\$readmemh} and \code{\$readmemb} have no \VHDL{} equivalent short of a
TEXTIO loop you write yourself, and they are worth knowing well: one call
loads an entire unpacked array from a text file of hexadecimal or binary
values, they accept \code{@address} lines to jump around, and they are
supported by every synthesis tool as a way to initialise a ROM.

\begin{svcode}[caption={A ROM initialised from a file.},label={lst:rosetta:rom}]
module atan_rom #(parameter int unsigned N = 16) (
  input  wire  [3:0]  addr,
  output logic [15:0] data
);
  logic [15:0] rom [0:N-1];
  initial $readmemh("atan_table.hex", rom);
  assign data = rom[addr];
endmodule
\end{svcode}

% ===========================================================================
\section{Time and delays}
\label{sec:rosetta:time}
% ===========================================================================

\begin{rosetta}
\begin{rleft}
-- time is a physical type with
-- units built into the language
constant PERIOD : time := 10 ns;

wait for PERIOD;
wait for 100 ps;
wait for 1 us;

-- resolution limit is set by the
-- simulator command line, e.g.
--   vsim -t ps work.tb

-- now / current time:
report time'image(now);
\end{rleft}
\begin{rright}
`timescale 1ns/1ps
// or, inside the module:
timeunit      1ns;
timeprecision 1ps;

localparam time PERIOD = 10ns;

initial begin
  clk = 0;  #PERIOD;
  clk = 1;  #100ps;
  clk = 0;  #1us;
  clk = 1;  #10;
  // "#10" is 10 TIME UNITS, and
  // which unit that is depends
  // on the timescale in force.
  $display("%0t", $time);
end
\end{rright}
\end{rosetta}

\index{timescale directive}\index{time!units}

\VHDL{}'s \code{time} is a physical type: a value carries its unit, the
compiler checks it, and \code{10 ns} means ten nanoseconds no matter where it
appears. The simulator's resolution limit affects rounding, not meaning.

\SV{} inherited Verilog's model, in which a delay is a bare number scaled by a
directive. \code{\#10} means ``ten time units'', and which unit that is depends on
the \chTwoTick{timescale} in force for the module. \SV{} added time literals
with units (\code{\#10ns}) and the \kw{timeunit}/\kw{timeprecision} declarations,
and between them they make the problem avoidable. It is still there for anyone
who does not use them.

\begin{gotcha}{The timescale rule is genuinely awful}
The timescale that applies to a module is the one established by the last
\chTwoTick{timescale} directive that the compiler read \emph{before} that
module, in the order the files were given on the command line. Not the file's
own directive, if it has none; not the project's; the last one seen. So:

\begin{itemize}
  \item A module in a file with no \chTwoTick{timescale} inherits whatever the
        previous file set. Change the file order and its delays change.
  \item If no file sets one, the timescale is tool-defined, and
        \tool{QuestaSim}, \tool{VCS} and \tool{Xcelium} do not agree.
  \item The simulation's time precision is the \emph{smallest} precision of any
        module in the design, so one testbench file with \code{1ps} precision
        slows down the whole simulation.
  \item \code{\$time} returns the time scaled and rounded to the calling
        module's time unit, so the same instant prints differently in two
        modules. \code{\$realtime} avoids the rounding but not the scaling.
\end{itemize}

\begin{svcode}
// file a.sv
`timescale 1ns/1ps
module a; initial #10 $display("a at %0t", $time); endmodule
// file b.sv -- no directive of its own
module b; initial #10 $display("b at %0t", $time); endmodule
\end{svcode}
Compile \code{a.sv b.sv} and \code{b} waits 10\,ns. Compile \code{b.sv a.sv}
and \code{b}'s timescale is tool-defined. Nothing warns you.

The fix is mechanical and you should apply it to every file you write:
\begin{svcode}
`timescale 1ns/1ps        // first line of every .sv file
module m;
  timeunit      1ns;      // belt and braces: module-local, order-independent
  timeprecision 1ps;
  ...
endmodule
\end{svcode}
and write delays with explicit units, \code{\#10ns}, never bare \code{\#10}.
\end{gotcha}

\begin{notebox}{Delays do not belong in \RTL{} at all}
Both languages let you write a delay in synthesisable code and both ignore it
at synthesis. The result is a simulation that does not match the netlist. The
only delays in a well-written project are in the clock generator and the
stimulus of the testbench:
\begin{svcode}
localparam time TCLK = 10ns;
initial begin
  clk = 1'b0;
  forever #(TCLK/2) clk = ~clk;
end
\end{svcode}
Everything else is driven from clock edges.
\end{notebox}

% ===========================================================================
\section{The complete mapping table}
\label{sec:rosetta:table}
% ===========================================================================

\cref{tab:rosetta:master} is the summary. It repeats, in one place, everything
that \cref{ch:rosetta} and this chapter have established, static half and
behavioural half together. Photocopy it.

\begin{longtable}{@{}>{\raggedright\arraybackslash}p{0.31\linewidth}>{\raggedright\arraybackslash}p{0.28\linewidth}>{\raggedright\arraybackslash}p{0.35\linewidth}@{}}
  \caption{\VHDL{} to \SV{}: the complete mapping.}
  \label{tab:rosetta:master} \\
  \toprule
  \textbf{\VHDL{}} & \textbf{\SV{}} & \textbf{Notes} \\
  \midrule
  \endfirsthead
  \multicolumn{3}{@{}l}{\small\itshape \cref{tab:rosetta:master}, continued}\\
  \toprule
  \textbf{\VHDL{}} & \textbf{\SV{}} & \textbf{Notes} \\
  \midrule
  \endhead
  \midrule
  \multicolumn{3}{r@{}}{\small\itshape continues on the next page}\\
  \endfoot
  \bottomrule
  \endlastfoot

  \multicolumn{3}{@{}l}{\textit{Design units}}\\
  \midrule
  \code{entity} + \code{architecture} & \code{module} ... \code{endmodule} & One unit, not two. \\
  several architectures & parameter, \chTwoTick{ifdef}, or several modules & No ``same interface, other body''. \\
  \code{configuration} & \code{config} (rare) & Effectively no equivalent. \\
  \code{package} + \code{package body} & \code{package} ... \code{endpackage} & One unit. No deferred constants. \\
  \code{library ieee;} & --- & No library clause. \\
  \code{use pkg.all} & \code{import pkg::*} & Different resolution rules. \\
  \code{use pkg.name} & \code{import pkg::name} & \\
  \code{pkg.name} & \code{pkg::name} & Qualified reference, no import needed. \\
  --- & \chTwoTick{include} \file{"x.svh"} & Textual inclusion. Macros only. \\
  --- & \chTwoTick{define}, \chTwoTick{ifdef} & Preprocessor. No \VHDL{} equivalent. \\
  --- & \chTwoTick{default\_nettype none} & Put it in every file. \\
  analysis order & compilation unit, \code{\$unit} & Tool-defined. Avoid \code{\$unit}. \\
  \midrule

  \multicolumn{3}{@{}l}{\textit{Ports and parameters}}\\
  \midrule
  \code{in} / \code{out} / \code{inout} & \code{input} / \code{output} / \code{inout} & \\
  \code{buffer} & \code{output} & \SV{} outputs are always readable. \\
  port default \code{:= '0'} & --- & No default port values. \\
  \code{generic} & \code{parameter} & \\
  architecture \code{constant} & \code{localparam} & \\
  \code{generic (type T)} (2008) & \code{parameter type T} & Well supported in \SV{}. \\
  \code{generic map (X => 1)} & \code{\#(.X(1))} & Before the instance name. \\
  --- & \code{defparam} & Deprecated. Never use. \\
  \midrule

  \multicolumn{3}{@{}l}{\textit{Types}}\\
  \midrule
  \code{std\_ulogic} & \code{logic} & 4-state variable. \\
  \code{std\_logic} & \code{wire} & 4-state resolved net. \\
  \code{std\_logic\_vector(N-1 downto 0)} & \code{logic [N-1:0]} & \\
  \code{unsigned} & \code{logic [N-1:0]} & Unsigned is the default. \\
  \code{signed} & \code{logic signed [N-1:0]} & \\
  \code{bit} & \code{bit} & 2-state in both. \\
  \code{boolean} & --- & Any non-zero value is true. \\
  \code{integer} & \code{int} & Not \SV{}'s \code{integer}, which is 4-state. \\
  \code{natural}, \code{positive} & --- & No range constraints, no run-time checks. \\
  \code{real} & \code{real} & \\
  \code{character} & \code{byte} & \\
  \code{string} & \code{string} & Dynamic in \SV{}. \\
  \code{time} & \code{time} + \chTwoTick{timescale} & See the timescale trap. \\
  \code{subtype ... range} & --- & No equivalent. Assert instead. \\
  \code{type ... is array} & unpacked \code{[..]} after the name & For memories. \\
  --- & packed \code{[..]} before the name & One vector with a shape. \\
  \code{record} & \code{struct packed} / \code{struct} & Packed gives a defined bit layout. \\
  --- & \code{union}, \code{union packed} & No \VHDL{} equivalent. \\
  \code{type ... is (A, B)} & \code{typedef enum logic [..] \{A,B\}} & Always give the base type. \\
  \midrule

  \multicolumn{3}{@{}l}{\textit{Literals and conversions}}\\
  \midrule
  \code{'0'}, \code{'1'} & \code{1'b0}, \code{1'b1} & \\
  \code{"1010"} & \code{4'b1010} & \\
  \code{x"A5"} & \code{8'hA5} & \\
  \code{16\#FF\#} & \code{255} or \code{8'hFF} & \\
  \code{(others => '0')} & \code{'0} & Not \code{1'b0}. \\
  \code{(others => '1')} & \code{'1} & \\
  aggregate \code{(a => 1, b => 2)} & \code{'\{a:1, b:2\}} & Note the apostrophe. \\
  \code{\&} (concatenation) & \code{\{ , \}} & \code{\&} is bitwise AND in \SV{}. \\
  --- & \code{\{N\{x\}\}} & Replication. No \VHDL{} form. \\
  \code{unsigned(v)}, \code{std\_logic\_vector(u)} & --- & No conversion needed. \\
  \code{to\_integer(u)} & implicit, or \code{int'(u)} & \\
  \code{to\_unsigned(i,N)} & implicit truncation & \\
  \code{to\_signed(i,N)} & \code{\$signed(...)} & \\
  \code{resize(v,N)} & implicit, or \code{N'(v)} & Silent. This is the danger. \\
  \code{integer(r)} on a real & implicit, rounds & \\
  --- & \code{\$cast(d,s)} & Checked conversion at run time. \\
  --- & \code{\{>>\{x\}\}} & Bit-stream cast. \\
  \midrule

  \multicolumn{3}{@{}l}{\textit{Statements}}\\
  \midrule
  \code{y <= expr;} (concurrent) & \code{assign y = expr;} & \\
  \code{y <= a when c else b;} & \code{assign y = c ? a : b;} & \\
  \code{with s select y <= ...} & \code{always\_comb} + \code{case} & No concurrent form. \\
  \code{process (a, b)} & \code{always @(a, b)} & \\
  \code{process (all)} & \code{always\_comb} & Also runs at time zero. \\
  \code{process (clk)} + \code{rising\_edge} & \code{always\_ff @(posedge clk)} & \\
  \code{process (clk, rst\_n)} & \code{always\_ff @(posedge clk or negedge rst\_n)} & Asynchronous reset. \\
  signal \code{<=} in a process & non-blocking \code{<=} & Use in \code{always\_ff}. \\
  variable \code{:=} in a process & blocking \code{=} & Use in \code{always\_comb}. \\
  --- & \code{initial} & Process with no sensitivity list. \\
  --- & \code{final} & No \VHDL{} equivalent. \\
  \code{if / elsif / else / end if} & \code{if / else if / else} & \code{begin}/\code{end} required for blocks. \\
  \code{case} + \code{when others} & \code{case} + \code{default} & \SV{} does not require completeness. \\
  \code{case?} (2008) & \code{casez} & Never \code{casex}. \\
  --- & \code{unique}, \code{unique0}, \code{priority} & Full/parallel case, plus a run-time check. \\
  \code{for i in r loop} & \code{for (int i=..)} / \code{foreach} & \\
  \code{while}, \code{loop} & \code{while}, \code{forever}, \code{repeat}, \code{do..while} & \\
  \code{next}, \code{exit} & \code{continue}, \code{break} & \\
  \code{null} & \code{;} & \\
  \code{wait for 10 ns} & \code{\#10ns} & \\
  \code{wait until c = '1'} & \code{wait (c)} & Different: no edge required. \\
  \code{wait on a, b} & \code{@(a or b)} & \\
  \code{after}, \code{transport} & \code{\#} & No inertial/transport distinction. \\
  \midrule

  \multicolumn{3}{@{}l}{\textit{Structure}}\\
  \midrule
  \code{for ... generate} & \code{for (genvar i=..) begin : g\_x} & Name every block. \\
  \code{if ... generate} & \code{if (..) begin : g\_x} & \\
  \code{case ... generate} (2008) & \code{case (..) ... endcase} & \\
  \code{block} & named \code{begin}/\code{end} or generate block & \\
  \code{component} + \code{port map} & direct instantiation & No component declaration. \\
  \code{entity work.f(rtl)} & \code{f u\_f (...)} & No architecture selection. \\
  \code{=> open} & \code{.p()} & Leave nothing implicit. \\
  --- & \code{.name}, \code{.*} & Implicit connection. Use \code{.*} sparingly. \\
  --- & \code{bind} & Attach checkers from outside. No \VHDL{} form. \\
  \midrule

  \multicolumn{3}{@{}l}{\textit{Subprograms and reporting}}\\
  \midrule
  \code{function} & \code{function automatic} & Always write \code{automatic}. \\
  \code{procedure} & \code{task automatic} & Or \code{function void}. \\
  \code{signal} parameter & \code{ref} & \\
  \code{impure} & (default) & No purity control in \SV{}. \\
  overloading, operator overloading & --- & No equivalent. \\
  \code{assert ... report ... severity error} & \code{assert (c) else \$error(..)} & Fires when false in both. \\
  \code{note}, \code{warning}, \code{error}, \code{failure} & \code{\$info}, \code{\$warning}, \code{\$error}, \code{\$fatal} & \\
  \code{report x \& integer'image(i)} & \code{\$display("x \%0d", i)} & \\
  \code{std.textio} & \code{\$fopen}, \code{\$fscanf}, \code{\$fdisplay} & No \code{line} object. \\
  --- & \code{\$readmemh}, \code{\$readmemb} & Memory initialisation. \\
\end{longtable}

% ===========================================================================
\section{Checklist}
\label{sec:rosetta2:checklist}
% ===========================================================================

Operational takeaways for the behavioural half. When you see the left, write
the right, and never do the things marked ``never''. The static half is at the
end of \cref{ch:rosetta} and is not repeated here, with one exception: keep
running \tool{Verilator} with \code{-Wall} over every file and treating
\code{WIDTHTRUNC}, \code{WIDTHEXPAND} and \code{PINMISSING} as errors, because
none of what follows protects you from a width.

\begin{itemize}
  \item \textbf{Every module repeats its timescale} as
        \kw{timeunit}/\kw{timeprecision}, because
        \chTwoTick{timescale} is inherited across files in command-line order
        and nothing warns you when it is wrong.
  \item \textbf{Processes:} \code{process (all)} becomes \code{always\_comb};
        \code{process (clk)} with \code{rising\_edge} becomes
        \code{always\_ff @(posedge clk)}. Use blocking \code{=} in
        \code{always\_comb} and non-blocking \code{<=} in \code{always\_ff},
        never the other way round, never both in one block.
  \item \textbf{One writer per variable.} \code{always\_comb} and
        \code{always\_ff} enforce this; plain \code{always} and \code{initial}
        do not, and module-level variables are shared by every process.
  \item \textbf{\code{case}:} always write a \code{default} in combinational
        logic, because \SV{} does not require completeness and the missing
        branch is a latch. Prefer \code{unique case} to get a run-time check
        as well. Never write \code{casex}.
  \item \textbf{Name every generate block} with a \code{g\_} prefix, or the
        tool numbers them and your hierarchical paths change when you edit the
        file.
  \item \textbf{Instantiate with full named association} in \RTL{}. There is no
        component declaration to keep in step, but there is also no width
        check: \code{.*} belongs in testbench wrappers only.
  \item \textbf{\code{wait until c = '1'} is \emph{not} \code{wait (c)}.}
        \VHDL{} waits for an event that makes the condition true; \SV{} returns
        at once if it already holds. Where you meant the \VHDL{} behaviour,
        write \code{@(posedge c)}.
  \item \textbf{Assertions:} \code{assert ... report ... severity} becomes
        \code{assert (c) else \$error(...)}, so the message moves into the
        \code{else}. Both fire when the condition is false.
  \item \textbf{Every function and task is \kw{automatic}.} A static
        subprogram breaks recursion and races when two processes call it.
  \item \textbf{Put types, constants and functions in a package}, macros in an
        \file{.svh} header, and nothing at all in \code{\$unit}.
  \item \textbf{Reach for \code{\$clog2}, \code{\$countones}, \code{\$onehot},
        \code{\$isunknown} and \code{\$bits}} where \VHDL{} made you write a
        function. Do not expect \code{'range}, \code{'reverse\_range},
        \code{'delayed} or \code{'value}: they have no equivalent.
  \item \textbf{Delays belong in testbenches only}, written with explicit units
        (\code{\#10ns}), and everything else is driven from a clock edge.
\end{itemize}

Together with \cref{ch:rosetta}, this chapter is enough to read and write
\SV{} \RTL{}: every construct above has a \VHDL{} original you already know.
What comes next has no original. \cref{ch:beyond} takes the design-side
constructs \SV{} has and \VHDL{} does not, starting with interfaces, packed
structures and parameterised types, and \refverificationbook takes the
verification-side ones, from classes and constrained randomisation to the
assertion sublanguage. Those two chapters are what makes the move worth
making.
